% !TEX root = ../diary-book.tex



%
%\section{令和7年10月17日 中間報告にむけての準備}
%
%論文を書くまではいかなくとも、どの程度の精度で測位が可能かを示すことができるとよい。
%また、応用先についての具体的な検討が必要であるので、それを含めて、どの程度までいけるかを整理する。
%その他の測位手法との差異についても整理を行う必要がある。
%
%\section{令和7年10月22日 ユースケースの一つとしての潜水艦の測位}
%
%これまでにいただいた情報を整理すると、潜水艦の測位のユースケースとしては以下のような条件が想定される。
%
%\begin{itemize}
%  \item 深さ：～500m
%  \item 速度：6km/h
%  \item 範囲：1km四方
%\end{itemize} 
%\subsection{深さとミューオンの到達可能性}
%深さ 500 m の場合、海水の密度は 1.03 g/cm$^3$ 程度であるので、500 m の水深では 515 m の水の厚みがあることになる。空気中の大気圧は 1 atm = 1013 hPa = 1013 mbar = 1013 g/cm$^2$ であるので、500 m の水深では 515 g/cm$^2$ の厚みがあることになる。ミューオンのエネルギー損失は 1 g/cm$^2$ あたり 2 MeV 程度であるので、500 m の水深では 1 GeV 程度のエネルギー損失があることになる。したがって、500 m の水深においても、1 GeV 以上のエネルギーを持つミューオンは到達することができる。ただし、1~GeV はギリギリのエネルギーであるので、2 GeV 以上のエネルギーを持つミューオンが到達することができると考えたほうがよい。また、測位には角度がついたミューオンが必要である。４５度程度までのミューオンを考えるとするとエネルギー損失は 1.4~GeV であり、2~GeV 以上のエネルギーを持つミューオンの到達が考えられる。
%
%\begin{codebox}{連絡先}
%  \begin{minted}{text}
%***************************************
%株式会社IHI
%航空・宇宙・防衛事業領域　防衛事業連携統括部
%高橋　慶貴
%TEL：03-6204-7120
%MAIL：takahashi0046@ihi-g.com
%***************************************
%\end{minted}
%\end{codebox}
%
%
%
